% !TEX root = ../main.tex
%% Sujet type bac n° 2 — sujet inédit OptiTAB (calculs vérifiés par programme : bac_s2.py)

\sujetBac{2}{QCM justifié : analyse et probabilités/4,Suite arithmético-géométrique/5,Géométrie dans l'espace/4,Fonction logarithme et intégration/5}

\begin{exoBac}{4}{Questionnaire à choix multiples justifié}
Pour chacune des questions suivantes, une seule des quatre réponses proposées est exacte. Indiquer sur la copie le numéro de la
question et la lettre de la réponse choisie, puis \textbf{justifier ce choix} par un calcul ou un raisonnement rédigé. Une réponse
non justifiée ne rapporte aucun point.
\begin{questions}
  \item La limite en $+\infty$ de la fonction $x\mapsto\dfrac{e^{x}-x^{2}}{2e^{x}+1}$ est égale à :
\begin{tasks}[label=\alph*)](4)
  \task $0$
  \task $\frac12$
  \task $1$
  \task $+\infty$
\end{tasks}
  \item Une primitive sur $\R$ de la fonction $f:x\mapsto x\,e^{x^{2}}$ est la fonction $F$ définie par :
\begin{tasks}[label=\alph*)](4)
  \task $F(x)=x^{2}e^{x^{2}}$
  \task $F(x)=2e^{x^{2}}$
  \task $F(x)=e^{x^{2}}$
  \task $F(x)=\frac12e^{x^{2}}$
\end{tasks}
  \item Sur l'intervalle $\intOO1{+\infty}$, l'ensemble des solutions de l'équation $\ln x+\ln(x-1)=\ln6$ est :
\begin{tasks}[label=\alph*)](4)
  \task $\{3\}$
  \task $\{-2\pv3\}$
  \task $\{2\}$
  \task l'ensemble vide
\end{tasks}
  \item Une main est un ensemble de $5$ cartes choisies dans un jeu de $32$ cartes, qui contient $4$ as. Le nombre de mains contenant
    exactement deux as est :
\begin{tasks}[label=\alph*)](4)
  \task $\binom42+\binom{28}{3}$
  \task $4^{2}\times28^{3}$
  \task $\binom42\times\binom{28}{3}$
  \task $\binom{32}{5}-\binom{28}{5}$
\end{tasks}
  \item La variable aléatoire $X$ suit la loi binomiale $\mathcal{B}(10\pv0{,}3)$. Une valeur approchée au millième de $P(X\geq1)$ est :
\begin{tasks}[label=\alph*)](4)
  \task $0{,}028$
  \task $0{,}121$
  \task $0{,}850$
  \task $0{,}972$
\end{tasks}
\end{questions}
\end{exoBac}

\begin{exoBac}{5}{Suite arithmético-géométrique}
Au 1\textsuperscript{er} janvier 2026, une lettre d'information compte $2\,000$ abonnés. On estime que, chaque mois, $10\,\%$ des
abonnés se désabonnent et que $300$ nouvelles personnes s'abonnent. On modélise le nombre d'abonnés $n$ mois après le
1\textsuperscript{er} janvier 2026 par une suite $(u_n)$, avec $u_0=2\,000$.
\begin{questions}
  \item Calculer $u_1$ et $u_2$.
  \item Justifier que, pour tout entier naturel $n$, $u_{n+1}=0{,}9\,u_n+300$.
  \item On pose, pour tout entier naturel $n$, $v_n=u_n-3\,000$.
    \begin{enumerate}[label=\alph*),leftmargin=1.4em,itemsep=0.4mm,topsep=0.4mm]
      \item Démontrer que la suite $(v_n)$ est géométrique ; préciser sa raison et son premier terme.
      \item En déduire que, pour tout entier naturel $n$, $u_n=3\,000-1\,000\times0{,}9^{n}$.
    \end{enumerate}
  \item Étudier le sens de variation de la suite $(u_n)$, puis déterminer sa limite. Interpréter ces résultats.
  \item Le responsable souhaite dépasser $2\,900$ abonnés.
    \begin{enumerate}[label=\alph*),leftmargin=1.4em,itemsep=0.4mm,topsep=0.4mm]
      \item Résoudre dans l'ensemble des entiers naturels l'inéquation $u_n>2\,900$. Au cours de quel mois l'objectif sera-t-il atteint ?
      \item Recopier et compléter la fonction Python ci-dessous pour qu'elle renvoie la plus petite valeur de $n$ telle que $u_n>s$.
    \end{enumerate}
\begin{python}
def seuil(s):
    u = 2000
    n = 0
    while ... :
        u = ...
        n = n + 1
    return n
\end{python}
\end{questions}
\end{exoBac}

\begin{exoBac}{4}{Géométrie dans l'espace}
L'espace est muni d'un repère orthonormé. On considère les plans $\mathcal{P}_1:x+y-z-1=0$ et $\mathcal{P}_2:2x-y+z-2=0$, ainsi que le
point $M(0\pv1\pv3)$.
\begin{questions}
  \item Justifier que les plans $\mathcal{P}_1$ et $\mathcal{P}_2$ sont sécants.
  \item Montrer que leur droite d'intersection $\Delta$ admet pour représentation paramétrique
    \[\begin{cases}x=1\\y=t\\z=t\end{cases}\qquad t\in\R.\]
  \item Les plans $\mathcal{P}_1$ et $\mathcal{P}_2$ sont-ils perpendiculaires ?
  \item Déterminer une équation cartésienne du plan $\mathcal{Q}$ passant par $M$ et orthogonal à la droite $\Delta$.
  \item Déterminer les coordonnées du point $H$, intersection de la droite $\Delta$ et du plan $\mathcal{Q}$.
  \item Justifier que $H$ est le projeté orthogonal de $M$ sur $\Delta$, puis calculer la distance du point $M$ à la droite $\Delta$.
  \item Déterminer les coordonnées du point $M'$, symétrique de $M$ par rapport à la droite $\Delta$.
\end{questions}
\end{exoBac}

\begin{exoBac}{5}{Fonction logarithme et intégration}
On considère la fonction $f$ définie sur $\intOO0{+\infty}$ par $f(x)=x^{2}-2\ln x$, et l'on note $\mathcal{C}_f$ sa courbe représentative.
\partieExo{A}{Étude de la fonction}
\begin{questions}
  \item Déterminer les limites de $f$ en $0$ et en $+\infty$.
  \item Montrer que $f'(x)=\dfrac{2(x-1)(x+1)}{x}$ et dresser le tableau de variations de $f$.
  \item En déduire que, pour tout réel $x>0$, $f(x)\geq1$.
  \item Étudier la convexité de $f$. Que peut-on en déduire pour la position de $\mathcal{C}_f$ par rapport à sa tangente au point
    d'abscisse $1$ ?
\end{questions}
\partieExo{B}{Équation et aire}
\begin{questions}[resume]
  \item Montrer que l'équation $f(x)=2$ admet exactement deux solutions $a$ et $b$, avec $a<1<b$. Donner un encadrement de $b$
    d'amplitude $10^{-2}$.
  \item On admet que la fonction $F:x\mapsto x\ln x-x$ est une primitive de la fonction $\ln$ sur $\intOO0{+\infty}$. Calculer
    $\int_1^{e}f(x)\,\mathrm{d}x$, puis interpréter ce résultat en termes d'aire.
\end{questions}
\end{exoBac}
